\paragraph{An interaction-generated charge scale in the original one-dimensional law.}
There is a positive many-body charge threshold without inserting a Dirac
mass.  On the unchanged periodic Weyl/contact geometry $(L,1,1)$,
$L=4m$, the onsite frame $G_x=(-1)^x(i\sigma_x)^x$ satisfies
$G_x(i\sigma_x)G_{x+1}^\dagger=-I$ and has period four.  Its CAR lift
fixes local doublons, so the original action is exactly the repulsive
Hubbard chain with $t=c/2=1/16$ and $U=1$.  Four exact local frame
identities, source-matrix checks through $L=16$, and an independent
256-state full-Fock witness bind this join.  Unlike the earlier
logarithmic-reference alternative Hamiltonian, no coupling is reassigned.

For the \emph{half-filled canonical ground sector}, the imported
\href{https://arxiv.org/html/cond-mat/0207529}{Lieb--Wu charge-gap theorem}
gives
\[
 \Delta_c=\lim_{L\to\infty}
 [E_L(L+1)+E_L(L-1)-2E_L(L)]
 =U-4t+8t\sum_{k\ge1}{1/2\choose k}
       (2t/U)^{2k-1}\eta(2k-1).
\]
Here the series is used only for $2t/U<1$.  Twelve binomial terms and
96-term rational Euler bounds for each Dirichlet eta value enclose
$\Delta_c=0.7715517157380375\ldots$ with width below $1.67\,10^{-25}$.
The binomial tail is bounded by
$8t|{1/2\choose K+1}|(2t/U)^{2K+1}/[1-(2t/U)^2]$;
the nonnegative Euler coefficients give tail $2^{-96}$.
Floating quadrature is an independent check only.

Particle--hole symmetry fixes the charge costs for the original $H$:
\[
 E_+=\frac{U+\Delta_c}{2}=0.885775857869\ldots,\qquad
 E_-=\frac{\Delta_c-U}{2}=-0.114224142131\ldots .
\]
Thus half filling is a specified canonical reference, not the original
$\mu=0$ grand-canonical vacuum.  The symmetric threshold
$\Delta_c/2=0.3857758578690188\ldots$ uses the explicitly centered
$K=H-(U/2)N$; $\Delta_c$ itself is invariant under any common $-\mu N$.
The spin sector remains gapless, contact vanishes in the
one-particle sector, and no preparation or complete-repair preservation
of this many-body reference is established by the gap theorem.

\paragraph{The fixed-interaction continuum boundary.}
The equivalent positive integral
\href{https://arxiv.org/html/cond-mat/9911281}{for the charge gap} and
the elementary $K_1$ integral bounds imply, for every $t,U>0$,
\[
 \Delta_0\le\Delta_c\le
 \Delta_0\frac{1+3/(8A)}{1-e^{-2A}},\qquad
 A=\frac{2\pi t}{U},\qquad
 \Delta_0=\frac8\pi\sqrt{Ut}\,e^{-A}.
\]
Consequently fixed $U$ with $t=v/(2a)$ has
$\Delta_c\to0$ as $a\to0$: this lattice gap does not itself supply a
nonzero continuum mass.  A sufficient additional running family is
\[
 a(A)=\frac{4v e^{-A}}{M\sqrt{2\pi A}},\qquad
 t(A)=\frac{v}{2a(A)},\qquad U(A)=\frac{\pi v}{a(A)A},
\]
for a \emph{chosen} $M,v>0$.  It gives $\Delta_0/2=M$ exactly and
$1\le\Delta_c/(2M)\le[1+3/(8A)]/[1-e^{-2A}]\to1$.
This is a quantitative charge-threshold limit, not a derivation of $M$
from the fixed source, a complete relativistic scaling theory, or a
$3+1$-dimensional massive species.  The receipt
\texttt{model/mass\_generation\_audit/results.json} keeps these reference,
coupling-family, and dimensional premises explicit.
